% ==============================================================================
% dxh_flashing.tex — DXH workshop Part 2: dev environment + firmware flashing
% DXH講座 ②: 開発環境 + 書き込み体験
% ==============================================================================

% ------------------------------------------------------------------
\begin{frame}[fragile]{② 開発環境の全体像}
  \begin{block}{今日は構築済み}
    \footnotesize 貸出 PC には ESP-IDF（組み込み開発環境）と sf CLI をセットアップ済み。\\
    実習は書き込みからスタートします
  \end{block}

  \begin{lstlisting}[style=sfbash, numbers=none]
git clone https://github.com/M5Fly-kanazawa/stampfly_ecosystem
cd stampfly_ecosystem
install.bat    ← ESP-IDF + sf CLI を自動インストール
setup_env.bat  ← 開発環境をアクティベート
  \end{lstlisting}

  {\footnotesize 上記は「自分の PC でも試したい」ときの手順（今日は実施不要）。「←」以降は説明なので入力しません}

  \vspace{0.2em}

  \begin{exampleblock}{このあと}
    \footnotesize GUI フラッシャ（デスクトップアプリ）のデモも行います
  \end{exampleblock}
\end{frame}

% ------------------------------------------------------------------
\begin{frame}{② 書き込みの仕組み --- 3つの方法}
  \begin{block}{ファームウェアとは}
    \footnotesize 機体を動かすプログラム一式のこと。「書き込み」とは、新しい\\
    ファームウェアを StampFly の中に転送する作業を指す
  \end{block}

  \vspace{0.3em}

  \begin{table}
    \centering\footnotesize
    \begin{tabular}{l>{\raggedright\arraybackslash}p{28mm}>{\raggedright\arraybackslash}p{28mm}>{\raggedright\arraybackslash}p{34mm}}
      \hline
       & \textbf{Web フラッシャ} & \textbf{GUI フラッシャ} & \textbf{sf CLI（コマンド）} \\
      \hline
      起動方法     & ブラウザ（Edge/Chrome） & デスクトップアプリ & コマンドプロンプト \\
      書き込むもの & 公開されている完成版 & 公開されている完成版 & PC 内でビルドしたもの \\
      本日の用途   & 紹介のみ & デモのみ & \textbf{②・③の実習} \\
      \hline
    \end{tabular}
  \end{table}

  \vspace{0.3em}

  {\footnotesize 今日の実習は②・③とも sf CLI で書き込みます。Web / GUI フラッシャは\textbf{開発環境のない PC でも使える}方法（紹介のみ）。\\ ※ ビルド＝人が書いたコードを機体が実行できる形式に変換する作業}
\end{frame}

% ------------------------------------------------------------------
\begin{frame}[fragile]{② 実習：sf flash でファームウェアを書き込む}
  \begin{block}{PC 側の操作（コマンドプロンプト）}
    \begin{lstlisting}[style=sfbash, numbers=none]
cd stampfly_ecosystem
setup_env.bat
sf flash vehicle   ← 機体をUSB接続してから実行
    \end{lstlisting}
  \end{block}

  {\footnotesize ※ 「←」以降は説明なので入力しません}

  \vspace{0.2em}

  \begin{enumerate}\small
    \item 機体の電源を入れ、USB-C ケーブルで PC に接続
    \item 上のコマンドを実行（PC 内でビルド済みの標準ファームウェアを書き込み）
    \item 完了後、機体が再起動 --- 起動音が終わり LED が緑常灯になるまで水平静置（ジャイロ校正）
  \end{enumerate}
\end{frame}

% ------------------------------------------------------------------
\begin{frame}{② よくある失敗と対処}
  \begin{table}
    \centering\small
    \begin{tabular}{l>{\raggedright\arraybackslash}p{64mm}}
      \hline
      \textbf{症状} & \textbf{対処} \\
      \hline
      機体が認識されない/繋がらない & ケーブルを交換（充電専用ケーブルは不可） \\
      書き込みが進まない/失敗する & ポートを使う他アプリ（シリアルモニタ等）を閉じる \\
      それでも失敗する & 機体の BOOT ボタンを押しながら USB を挿し直して、もう一度 \code{sf flash} \\
      \hline
    \end{tabular}
  \end{table}

  \vspace{0.5em}

  {\footnotesize ※ \code{sf flash} は機体の設定（ペアリング等）には触れません。失敗したら、そのままもう一度実行するだけです\\
  ※ BOOT ボタン＝機体上面 M5StampS3 のボタン（③で ARM に使うものと同じ。起動時に押しながらだと書き込みモード）}
\end{frame}
